\documentclass[12pt]{article}
\usepackage[heading]{ctex}
\usepackage{array}
\usepackage{anyfontsize}
\usepackage{geometry}
\usepackage{graphicx}
\usepackage{fancyhdr}
\usepackage{ulem}
\usepackage{setspace}
\usepackage{lastpage}
\usepackage{fontspec}
\usepackage{fontawesome}
\usepackage{multirow}
\usepackage{multicol}
\usepackage{listings}
\usepackage{algorithm}
\usepackage{algorithmicx}
\usepackage{algpseudocode}
\usepackage{amsmath,amsthm,amsfonts,amssymb,bm}
% \usepackage{newtxtext,newtxmath}
\usepackage{epigraph}
\usepackage{hyperref}
\usepackage{tikz}
\usepackage{fancyvrb}
\usepackage{tcolorbox}
\usepackage{tabularx}
\usepackage{xeCJKfntef,xcolor}
%\usepackage{titlesec}
\usepackage{enumitem}
\usepackage{booktabs}
\hypersetup{
	colorlinks=true,
	linkcolor=blue,
	filecolor=blue,
	urlcolor=blue,
	citecolor=cyan,
}
\usetikzlibrary{arrows.meta}
\tikzset{global scale/.style={
	scale=#1,
	every node/.append style={scale=#1}
  }
}
\tcbset{
	colback=white,
	colframe=black,
	boxrule=0.5pt
}
\fvset{
	tabsize=0,
	frame=single,
	framesep=3mm,
	rulecolor=\color{blue!80},
	numbers=left,
	numbersep=5pt
}
\lstset{
	breaklines = true,
	keepspaces = true,
	basicstyle = \ttfamily,
	numberstyle = \footnotesize\ttfamily\color{gray},
	numbers = left,
	columns = flexible,
	framexleftmargin = 5pt,
	xleftmargin = 5pt,
	numbersep = 12pt,
	keywordstyle = \color{black}, %\color{blue!70},
	commentstyle = \color{black}, %\color{green!70},
	frame = single,
	frameround = tttt,
	rulecolor = \color{blue},
	language = c++,
	escapeinside = {(*@}{@*)}
}
\xeCJKsetup{
	underdot/symbol = {\color{black}.}
}
\ctexset{
	section = {
		format = \centering\Large\heiti,
		numbering = true,
		number = {},
		aftername = {}
	},
	subsection = {
		format = \fontsize{13pt}{\baselineskip}\heiti,
		numbering = true,
		number = {},
		aftername = {\quad【},
		aftertitle = {】}
	}
}
\geometry{
	a4paper,
	scale=0.8
}
\floatname{algorithm}{算法}
\renewcommand{\algorithmicrequire}{\textbf{输入：}}
\renewcommand{\algorithmicensure}{\textbf{输出：}}
\newcommand{\mrow}[2]{{\multirow{#1}{*}{#2}}}
\newcommand{\mcol}[2]{{\multicolumn{#1}{c|}{#2}}}
\newcommand{\CJKdotbf}[1]{{\textbf{\CJKunderdot{#1}}}}
\newcommand{\filename}[1]{{\textbf{\textit{#1}}}}
\newcommand{\slbf}[1]{\emph{\textbf{#1}}}
\newcommand{\Yes}[0]{$\sqrt{}$}
\newcommand{\No}[0]{$\times$}
\setCJKmainfont[BoldFont=SimHei,ItalicFont=KaiTi]{SimSun} % ------------------------------------------------- Windows
% \setCJKmainfont[BoldFont=FandolHei-Regular.otf,ItalicFont=FandolKai-Regular.otf]{FandolSong-Regular.otf} % -- Linux
\setmonofont{Consolas} % ------------------------------------------------------------------------------------ Windows
\setlist[itemize]{itemsep=0pt,partopsep=0pt,parsep=0pt,topsep=10pt}
\setlist[enumerate]{itemsep=0pt,partopsep=0pt,parsep=0pt,topsep=10pt}
\setlist[itemize,1]{leftmargin=35pt}
\setlist[enumerate,1]{leftmargin=35pt}
\newtheorem{lem}{引理}[section]
\newcommand{\docTitle}{NOIP模拟赛}
\newcommand{\docAuthor}{\text{Gold \& yyy}}
\newcommand{\fancyheadLeftText}{\docTitle}
\newcommand{\fancyheadRightText}{}
\title{\Huge \textbf{NOIP模拟赛}}
\author{$\docAuthor$}
\hypersetup{
	pdftitle={\docTitle},
	pdfauthor={\docAuthor}
}
\date{\heiti \Large{08:00$\sim$11:40}}
\begin{document}
\maketitle
\thispagestyle{fancy}
\pagestyle{fancy}
\renewcommand\headrulewidth{0pt}
\fancyhf{}
\fancyfoot[C]{\footnotesize 第 \thepage 页\quad 共 \pageref{LastPage} 页}
% \textwidth=478pt
% \tabcolsep=6pt
% total_width=\textwidth-2*columns*\tabcolsep-(columns+1)*0.5pt
\begin{center}
    \begin{tabular}{|m{60pt}|m{78.9020375pt}|m{78.9020375pt}|m{88.9020375pt}|m{108.9020375pt}|}
        \hline 题目名称 & 星光站 & 突触 & 虚空的召唤 & 十六号-塔 \\
        \hline 原题 & 无 & \href{https://codeforces.com/problemset/problem/1666/E}{Even} \href{https://www.luogu.com.cn/problem/P12818}{Split} & 无 & \href{https://atcoder.jp/contests/agc032/tasks/agc032_e}{Modulo Pairing} \\
        \hline
    \end{tabular}
\end{center}

\newpage
\renewcommand\headrulewidth{0.5pt}
\fancyhead[L]{\footnotesize \fancyheadLeftText}

\section{星光站（starlight）}
\fancyhead[R]{\footnotesize \fancyheadRightText\quad 星光站（starlight）}

题意转化为求有多少对 $(i,j)$ 使得 $i<j$ 且\textbf{不存在} $a_k$ 和 $x$ 使得 $a_i\oplus x<a_k\oplus x<a_j\oplus x$。

注意到如果存在 $a_j\oplus x<a_k\oplus x<a_i\oplus x$，那么将 $x$ 改为 $x\oplus (2^{31}-1)$ 则会出现 $a_i\oplus x<a_k\oplus x<a_j\oplus x$。

所以可以将条件改为 $i<j$ 且\textbf{不存在} $a_k$ 和 $x$ 使得 $a_i\oplus x<a_k\oplus x<a_j\oplus x$ 或者使得 $a_j\oplus x<a_k\oplus x<a_i\oplus x$。

于是可以发现 $i,j$ 是可以互换的，那么 $i<j$ 的条件就没用了，只需最后除以 $2$ 即可。

首先显然 $x=0$ 时要合法，所以先将 $a$ 从小到大排序，这样只有相邻两个会合法，所以考虑对每一个相邻对判断其是否合法。

设 $h(x,y)=\operatorname{highbit}(x\oplus y)$，即二进制下 $x,y$ 的最高不同位。

\begin{lem}
    若序列 $a$ 单调递增，则 $\forall 1\le i<j< |a|,h(a_i,a_j)\le h(a_i,a_{j+1})$。
\end{lem}

\begin{proof}
    由于 $a_i<a_j$，根据 $h(a_i,a_j)$ 表示二进制下 $a_i,a_j$ 的最高不同位可知 $a_i$ 这一位为 $0$，$a_j$ 这一位为 $1$。

    同理，$a_j$ 的第 $h(a_j,a_{j+1})$ 位为 $0$，所以 $h(a_i,a_j)\neq h(a_j,a_{j+1})$。

    显然，对于任意两个整数 $A, B$，若 $\operatorname{highbit}(A) \neq \operatorname{highbit}(B)$，则：
$$\operatorname{highbit}(A \oplus B) = \max(\operatorname{highbit}(A), \operatorname{highbit}(B))$$

    所以可以得到：

    $$
    \begin{aligned}
h(a_i, a_{j+1}) &= \operatorname{highbit}(a_i \oplus a_{j+1}) \\
&= \operatorname{highbit}\big((a_i \oplus a_j) \oplus (a_j \oplus a_{j+1})\big) \\
&= \max\big(\operatorname{highbit}(a_i \oplus a_j), \operatorname{highbit}(a_j \oplus a_{j+1})\big) \\
&= \max\big(h(a_i, a_j), h(a_j, a_{j+1})\big) \\
&\ge h(a_i, a_j)
    \end{aligned}
    $$

    
    
\end{proof}

\begin{lem}
    若 $a,b,c$ 互不相同且 $h(a,b)<h(b,c)$，则不存在非负整数 $x$ 使得 $a\oplus x<c\oplus x<b\oplus x$。
\end{lem}

\begin{proof}
    显然对于任意非负整数 $x$，均有 $h(a\oplus x,b\oplus x)=h(a,b)$，于是，只需证 $a<c<b$ 不成立。

    首先，$h(a,c)=\max\big(h(a,b),h(b,c)\big)=h(b,c)$。

    设 $H=h(a,c)=h(b,c)$，于是 $a,c$ 第 $H$ 位不同，$b,c$ 第 $H$ 位不同，所以 $a,b$ 第 $H$ 位相同，根据 $h()$ 的定义得若 $c$ 大于 $a$ 则 $c$ 大于 $b$，若 $c$ 小于 $a$ 则 $c$ 小于 $b$，所以不可能 $a<c<b$。
\end{proof}

于是，我们排完序之后就只需要判断 $h(a_i,a_{i+1})$ 是否同时小于 $h(a_{i+1},a_{i+2})$ 和 $h(a_{i-1},a_i)$ 即可。

时间复杂度 $O(n\log n)$，瓶颈在于排序。

\newpage

\section{突触（synapse）}
\fancyhead[R]{\footnotesize \fancyheadRightText\quad 突触（synapse）}

感觉要二分但是不会，那么先想想怎么判断是否有解。

不防考虑如何判断强制 $\min_{i}(x_i-x_{i-1})\ge L,\max_{i}(x_i-x_{i-1})\le R$ 时是否有解，即所有 $x_i-x_{i-1}$ 都在 $[L,R]$ 中时是否有解。

考虑设 $xl_i,xr_i$ 分别表示只考虑所有 $j\le i$，要求 $a_j\le x_j\le a_{j+1}$ 时，$x_i$ 能取 $[xl_i,xr_i]$，这个东西显然可以 $O(n)$ 递推：

\begin{itemize}
    \item $xl_i=\max(a_i,xl_{i-1}+L)$；
    \item $xr_i=\min(a_{i+1},xr_{i-1}+R)$。
\end{itemize}

随后只要判断 $l$ 是否在 $[xl_n,xr_n]$ 中即可。

发现上下两个式子一点关系都没有，所以可以将 $L,R$ 分别二分得到。

时间复杂度 $O(n\log V)$。

\newpage

\section{虚空的召唤（void）}
\fancyhead[R]{\footnotesize \fancyheadRightText\quad 虚空的召唤（void）}

首先我们可以将题目中 $N(u)$ 的定义改为与其直接相邻的点的集合，不包括自身，这显然不会影响答案。

如果我们确定了 $u$ 可以成为答案的其中一个点，那么我们显然可以 $O(n+m)$ 求出答案，只需将与 $u$ 相邻的点标记为 $1$，其余的点标记为 $0$，然后对每个点求出与其相邻的点中有多少个 $1$ 即可。

于是考虑找到答案中的一个点。

\begin{lem}
    如果 $u$ 的度数是奇数，那么 $u$ 一定可以成为答案的一个端点。
\end{lem}

\begin{proof}
    首先证明 $\sum_{v\in N(u)} |N(u)\cap N(v)|$ 为偶数。

    考察每一条边对这个的贡献，对于一条边 $\{x,y\}$，必须要当 $x$ 与 $u$ 直接相连，才会对 $|N(u)\cap N(y)|$ 贡献 $1$，但是只有当 $y$ 与 $u$ 也相邻时，才会对上面的式子贡献 $1$。但是此时，由于 $y$ 与 $u$ 相邻，$x$ 也与 $u$ 相邻，所以会再对上面的式子贡献 $1$。

    于是每条边要么对上面的式子贡献 $0$，要么贡献 $2$，所以 $\sum_{v\in N(u)} |N(u)\cap N(v)|$ 为偶数。

    由于与 $u$ 相邻的点数量为奇数，但是它们与 $u$ 的邻域交大小的和为偶数，奇数个数的和为偶数，所以其中一定有一个偶数。
\end{proof}

\begin{lem}
    如果所有点的度数都是偶数，那么对于任意一个点 $u$，$u$ 都可以成为答案的一个端点。
\end{lem}

\begin{proof}
    考虑 $\sum_{v\neq u}|N(u)\cap N(v)|$ 是什么。

    可以看成是先选一个不等于 $u$ 的点 $v$，再选一个与 $u,v$ 都相邻的点 $w$ 的方案数，那也就可以看成先选一个与 $u$ 直接相邻的点 $w$，再选一个与 $w$ 直接相邻的不等于 $u$ 的点 $v$ 的方案数，即 $\sum_{v\in N(u)}deg_v-deg_u$，可以发现这是个偶数。

    所以 $\sum_{v\neq u}|N(u)\cap N(v)|$ 为偶数，但是注意到总点数为偶数，去掉 $u$ 之后就只有奇数个点，那么奇数个数的和为偶数，其中就必有一个偶数。
\end{proof}

于是，我们可以 $O(n+m)$ 找出一个点，再 $O(n+m)$ 找出另一个点，时间复杂度 $O(n+m)$。

\newpage

\section{十六号-塔（tower）}
\fancyhead[R]{\footnotesize \fancyheadRightText\quad 十六号-塔（tower）}

接下来的 $n$ 代表题目中的序列长度。

题目相当于要让所有数两两匹配，要让匹配的两个数之和 $\bmod\ m$ 之后的最大值最小。

如果我们不 $\bmod\ m$，那这是一个相当简单的问题，直接将第一个数与最后一个数匹配，将第二个数与倒数第二个数匹配，以此类推即可。

如果 $\bmod\ m$ 会发生什么？如果两个数之和大于等于 $m$ 那么就会减去 $m$，所以不能直接贪心。

考虑将原序列复制成两份，左边一份都减 $m$，右边一份不变，然后拼接起来，得到序列 $b$。

发现任意一个长度为 $n$ 的区间都会包含原序列的所有数，现在我们只要给这个区间两两匹配使得和均非负，并且最大和最小，发现还是第一个数与最后一个数匹配等等，而且不难证明最优解一定是这其中一个区间的匹配方式。

于是我们只需要找到这个区间即可，也就是要找从两边往中间匹配后没有负数的最靠左的区间。

考虑依次找出长度为 $2,4,6,\dots n$ 的合法的最靠左区间。

初始设 $l=n+1,r=n$，每次 $l\leftarrow l-1,r\leftarrow r+1$，如果 $b_l+b_r<0$ 了就把两个数都加 $1$，如果 $b_{l-1}+b_{r-1}\ge 0$ 了就把两个数都减 $1$。

这样最后就能找到答案，时间复杂度 $O(n)$。

\end{document}
